\documentclass{article}
\usepackage{tcee-repaso}
\tcetema[Teoría de la demanda del consumidor (II)]{3.A.9}{Teoría de la demanda del consumidor (II). Dualidad e integrabilidad de las preferencias. Sistemas de demanda utilizados en estudios empíricos. Medidas de cambio en el bienestar.}
\tcefecha{14/10/2026}
% Dos gráficos lado a lado: \pareja{nombre1}{pie1}{nombre2}{pie2}
\newcommand{\pareja}[4]{\par\vspace{2pt}\noindent
  \begin{minipage}[b]{0.49\linewidth}\graficorepaso[\linewidth]{#1}{#2}\end{minipage}\hfill
  \begin{minipage}[b]{0.49\linewidth}\graficorepaso[\linewidth]{#3}{#4}\end{minipage}\par}
\begin{document}
\cabecerarepaso
\begin{columnas}

% ==========================================================================
\bloque{Esquema del cante}
\esquemacante{%
\apartadocante{Introducción}{2}
\apartadocante{I. Análisis teórico de la dualidad}{10}
\subcante{1. Ingredientes comunes}{2}
\subcante{2. Primal ($\vect{x}$, $V$) y dual ($\vect{h}$, $E$)}{5}
\subcante{3. Relaciones de dualidad}{3}
\apartadocante{II. Sistemas completos de demanda}{5}
\subcante{1. Idea}{0,5}
\subcante{2. Pasos: forma funcional (LES $\to$ AIDS), datos}{3,5}
\subcante{3. Valoración}{1}
\apartadocante{III. Aplicaciones de la dualidad}{12}
\subcante{1. Integrabilidad}{2,5}
\subcante{2. Ecuación de \autor{Slutsky}; clasificación}{3,5}
\subcante{3. Bienestar: $\Delta$EC, VC, VE, índices}{6}
\apartadocante{Conclusión}{1}}
Anexos (no se cantan): \autor{Kuhn}-\autor{Tucker} y CSO, envolvente (\autor{Roy}, \autor{Shephard}), hessiana, ordenación de medidas, \autor{Willig}.

% ==========================================================================
\bloque[2]{Introducción}
\begin{itemize}
\item \textbf{Enganche}: pensiones revalorizadas con el IPC (\dato{+2,7\,\% en 2026}{LGSS}): ¿basta para quedar igual? Si los precios suben de forma desigual, se sustituye: importa cuánto cuesta \emph{volver a estar igual de bien}. Utilidad ordinal $\Rightarrow$ hay que traducirla a euros: \concepto{función de gasto} $\Rightarrow$ \concepto{dualidad}.
\item \textbf{Relevancia}: \emph{teórica} (\autor{Slutsky} en dos líneas, integrabilidad); \emph{empírica} (elasticidades para IVA, IPC, pobreza); \emph{normativa} (VC y VE: coste-beneficio, exceso de gravamen).
\item \textbf{Contexto}: preferencias (no se ven) $\leftrightarrow$ demanda (se ve). Cuatro líneas: dualidad (\autor{Hotelling} 1932 $\to$ \autor{McKenzie} 1957, \autor{Diewert}); integrabilidad (\autor{Antonelli} 1886 $\to$ \autor{Hurwicz}-\autor{Uzawa} 1971); bienestar (\autor{Dupuit} 1844 $\to$ \autor{Hicks}, \autor{Willig}); sistemas (\autor{Stone} 1954 $\to$ AIDS 1980; Nobel a \autor{Deaton}, 2015).
\item \textbf{Problemática}: ¿qué se gana al «dar la vuelta» al problema? ¿Cómo se estiman las demandas? ¿Qué demandas son racionales? ¿Cómo medir el bienestar en euros y qué error cometen el excedente y los índices?
\item \textbf{Estructura}: I. dualidad (herramientas); II. sistemas de demanda (estimación); III. aplicaciones (integrabilidad, \autor{Slutsky}, bienestar).
\end{itemize}

% ==========================================================================
\bloque[10]{I. Análisis teórico de la dualidad}
\sub{I.1. Ingredientes comunes}
\begin{itemize}
\item $S\subseteq\mathbb{R}^n_+$ no vacío, cerrado, convexo. Axiomas: completitud, reflexividad, transitividad, continuidad ($\Rightarrow$ $U$ continua, \autor{Debreu}); insaciab.\ local; convexidad estricta; diferenciabilidad. Precios $\vect{p}\gg\vect{0}$.
\item Supuestos: estático, sin incertidumbre, info.\ perfecta, bienes homogéneos, precio-aceptante.
\end{itemize}


\sub{I.2. Primal y dual}
\begin{figure*}[t]
\centering\includegraphics[width=0.68\textwidth]{esquema-dualidad}\par
{\color{tceegris}\scriptsize\itshape Esquema de dualidad: problemas, soluciones, funciones valor y relaciones entre ellas (I.3)}
\end{figure*}
\begin{formula}
\[\text{PMU: }\max_{\vect{x}} U(\vect{x})\ \text{s.a.}\ \vect{p}\cdot\vect{x}\leq\overline{W}\ \Rightarrow\ \frac{U_1}{p_1}=\dots=\frac{U_n}{p_n}=\lambda\]
\[\text{PMG: }\min_{\vect{h}}\vect{p}\cdot\vect{h}\ \text{s.a.}\ U(\vect{h})\geq\overline{U}\ \Rightarrow\ \frac{p_1}{U_1}=\dots=\frac{p_n}{U_n}=\mu\]
\[\text{Ambos: } RMS_{12}=\frac{U_1}{U_2}=\frac{p_1}{p_2};\qquad \mu^*=\frac{1}{\lambda^*}\]
\end{formula}
\begin{itemize}
\item $\lambda$ = UMg de la renta $=\partial V/\partial\overline{W}$; $\mu$ = coste marginal de la utilidad $=\partial E/\partial\overline{U}$.
\item \concepto{Marshalliana} $\vect{x}(\vect{p},\overline{W})$: existe, única (\autor{Weierstrass}), continua (\autor{Berge}), homog.\ grado 0 en $(\vect{p},\overline{W})$ (sin ilusión monetaria), ley de \autor{Walras}.
\item \concepto{FIU} $V(\vect{p},\overline{W})=U(\vect{x}(\vect{p},\overline{W}))$: homog.\ 0, estr.\ creciente en $\overline{W}$, no creciente en $\vect{p}$, \textbf{cuasiconvexa} (no depende de la convexidad de las preferencias); identidad de \autor{Roy} (1947).
\end{itemize}
\pareja{primal}{Primal: la CI más alta alcanzable con $\overline{W}$}{dual}{Dual: la isogasto más baja que alcanza $\overline{U}$}
\begin{itemize}
\item \concepto{Hicksiana} $\vect{h}(\vect{p},\overline{U})$: homog.\ 0 en $\vect{p}$, sin exceso de utilidad, \concepto{ley de la demanda compensada} $\partial h_i/\partial p_i\leq0$ (no hay \autor{Giffen}: no tiene ER). \textbf{No observable} (depende de $\overline{U}$).
\item \concepto{Función de gasto} $E(\vect{p},\overline{U})=\vect{p}\cdot\vect{h}(\vect{p},\overline{U})$: homog.\ \textbf{grado 1} en $\vect{p}$, estr.\ creciente en $\overline{U}$, no decreciente en $\vect{p}$, \textbf{cóncava} en $\vect{p}$. Mide la utilidad en \emph{euros}.
\item Concavidad: $E(\vect{p}\,',\overline{U})\leq\vect{p}\,'\cdot\vect{h}(\vect{p},\overline{U})$: reoptimizar (sustituir) abarata respecto a la cesta fija.
\end{itemize}
\begin{formula}
\[\text{\autor{Roy}: } x_i=-\frac{\partial V/\partial p_i}{\partial V/\partial\overline{W}}\qquad \text{\autor{Shephard}: } h_i=\frac{\partial E}{\partial p_i}\]
\[\text{\autor{Slutsky}: } S_{ij}=\frac{\partial^2E}{\partial p_i\,\partial p_j}=\frac{\partial h_i}{\partial p_j}\]
\end{formula}
\begin{itemize}
\item $S$: \textbf{simétrica} (\autor{Schwarz}), \textbf{semidefinida negativa} ($E$ cóncava), $S\vect{p}=\vect{0}$ (\autor{Euler}).
\end{itemize}
\graficorepaso[0.46\linewidth]{concavidad-fg}{$E$ cóncava: siempre bajo su plano tangente, de pendiente $h_i$ (\autor{Shephard})}

\sub{I.3. Relaciones de dualidad}
\begin{itemize}
\item \concepto{Teorema básico} (MWG 3.E.1): con $U$ continua, insaciab.\ local y $\vect{p}\gg\vect{0}$, la cesta del PMU con $\overline{W}$ resuelve el PMG con $\overline{U}=U^*$ (gasto mínimo $=\overline{W}$), y la del PMG resuelve el PMU con $\overline{W}=E^*$ (utilidad máx.\ $=\overline{U}$).
\item Coinciden las \emph{cestas}, no las funciones: $\vect{x}$ y $\vect{h}$ dependen de variables distintas.
\end{itemize}
\graficorepaso[0.52\linewidth]{primal-dual}{Misma cesta: $\vect{x}^*=\vect{h}^*$ con $\overline{W}=E^*$, $\overline{U}=U^*$}
\begin{formula}
\[\vect{x}(\vect{p},\overline{W})\equiv\vect{h}\big(\vect{p},V(\vect{p},\overline{W})\big)\qquad \vect{h}(\vect{p},\overline{U})\equiv\vect{x}\big(\vect{p},E(\vect{p},\overline{U})\big)\]
\[E\big(\vect{p},V(\vect{p},\overline{W})\big)\equiv\overline{W}\qquad V\big(\vect{p},E(\vect{p},\overline{U})\big)\equiv\overline{U}\]
\end{formula}
\begin{itemize}
\item $V$ y $E$ \textbf{inversas} en su 2.º argumento. Basta resolver \emph{un} problema para tener las cuatro funciones (esquema de dualidad, arriba).
\item \emph{Cobb-Douglas} $U=x_1^{\alpha}x_2^{1-\alpha}$: $x_1=\alpha\overline{W}/p_1$, $V=\overline{W}\big(\frac{\alpha}{p_1}\big)^{\alpha}\big(\frac{1-\alpha}{p_2}\big)^{1-\alpha}$; invirtiendo, $E=\overline{U}\big(\frac{p_1}{\alpha}\big)^{\alpha}\big(\frac{p_2}{1-\alpha}\big)^{1-\alpha}$; \autor{Shephard}: $h_1=\alpha E/p_1$.
\end{itemize}
\begin{ideaclave}
Primal y dual son el mismo problema visto desde dos lados; $E$ mide la utilidad en dinero, sus derivadas son las hicksianas (\autor{Shephard}) y su concavidad da la ley de la demanda compensada.
\end{ideaclave}

% ==========================================================================
\bloque[5]{II. Sistemas completos de demanda}
\sub{II.1. Idea}
\begin{itemize}
\item Explican a la vez el gasto en \emph{todos} los bienes en función de precios y gasto total. Fines: \textbf{medir} elasticidades, \textbf{contrastar} la teoría (integrabilidad), \textbf{valorar} el bienestar.
\item Dualidad: se especifica $E$ o $V$ y se deriva (\autor{Shephard}/\autor{Roy}) $\Rightarrow$ coherente con la maximización \emph{por construcción}, sin resolver ningún problema.
\item \concepto{Requisitos} (tiran en sentidos opuestos): coherencia teórica, flexibilidad, agregación, sencillez (lineal).
\end{itemize}

\sub{II.2. Pasos}
\begin{itemize}
\item (1) forma funcional; (2) sistema completo; (3) datos; (4) econometría; (5) estimar y contrastar. Cada forma supera el límite de la anterior.
\item \concepto{Cobb-Douglas}: $s_i=\alpha_i$ constantes; $\varepsilon_W=1$, $\varepsilon_{ii}=-1$, cruzadas 0; homotética: no recoge \autor{Engel}.
\item \concepto{LES} (\autor{Stone}-\autor{Geary}, 1954): $U=\sum\alpha_i\ln(x_i-\gamma_i)$:
\end{itemize}
\begin{formula}
\[p_ix_i=p_i\gamma_i+\alpha_i\Big(\overline{W}-\sum_jp_j\gamma_j\Big)\]
\end{formula}
\begin{itemize}
\item Subsistencia + renta supernumeraria en proporciones fijas. \autor{Engel} rectas sin pasar por el origen. Límite: \emph{aditivo} (sin inferiores ni complementarios netos); \autor{Deaton} (1974) lo rechaza.
\item \concepto{Rotterdam} (\autor{Barten}, \autor{Theil}): \autor{Slutsky} en diferencias log.\ con coef.\ constantes; restricciones lineales; sin $E$ explícita.
\item \concepto{Translog} (\autor{Christensen}-\autor{Jorgenson}-\autor{Lau}, 1975): $\ln V$ de 2.º orden en $\ln(p_i/\overline{W})$, flexible; no lineal, no agrega, concavidad no garantizada.
\item \concepto{AIDS} (\autor{Deaton}-\autor{Muellbauer}, 1980): $E$ \textbf{PIGLOG} (\autor{Muellbauer}): $a(\vect{p})$ subsistencia, $b(\vect{p})$ saciedad.
\end{itemize}
\begin{formula}
\[\ln E=(1-\overline{U})\ln a(\vect{p})+\overline{U}\ln b(\vect{p})\]
\[s_i=\alpha_i+\sum_j\gamma_{ij}\ln p_j+\beta_i\ln\frac{\overline{W}}{P}\]
Agregación: $\sum_i\alpha_i=1$, $\sum_i\gamma_{ij}=\sum_i\beta_i=0$; homog.: $\sum_j\gamma_{ij}=0$; simetría: $\gamma_{ij}=\gamma_{ji}$.
\[\varepsilon_{x_i,\overline{W}}=1+\frac{\beta_i}{s_i}\]
\end{formula}
\begin{itemize}
\item $\beta_i>0$ lujo, $\beta_i<0$ necesario (\autor{Working}-\autor{Leser}). Agregación exacta (PIGLOG); \emph{LA-AIDS} con índice de \autor{Stone} $\ln P^*=\sum s_k\ln p_k$: MCO. Negatividad no garantizada (\emph{casi} ideal).
\item \concepto{QUAIDS} (\autor{Banks}-\autor{Blundell}-\autor{Lewbel}, 1997): término cuadrático en $\ln\overline{W}$, rango 3 (\autor{Gorman}): un bien puede ser de lujo para renta baja y necesario para renta alta.
\end{itemize}
\par\vspace{2pt}{\centering\renewcommand{\arraystretch}{1.1}\setlength{\tabcolsep}{2pt}\small%
\begin{tabular}{@{}>{\raggedright\arraybackslash}p{0.17\linewidth}>{\raggedright\arraybackslash}p{0.2\linewidth}>{\centering\arraybackslash}p{0.09\linewidth}>{\centering\arraybackslash}p{0.09\linewidth}>{\centering\arraybackslash}p{0.09\linewidth}>{\raggedright\arraybackslash}p{0.3\linewidth}@{}}
\toprule
\textit{Sistema} & \textit{Parte de} & \textit{Flex.} & \textit{Agreg.} & \textit{Lineal} & \textit{Límite}\\\midrule
C-D & $U$ directa & no & sí & sí (log) & $s_i$ constantes\\
LES & $U$ directa & no & sí & no & aditiva\\
Rotterdam & \autor{Slutsky} & local & no & sí & sin $E$ explícita\\
Translog & FIU & sí & no & no & no lineal\\
\rowcolor{tceesalvia}AIDS & $E$ PIGLOG & sí & exacta & LA & \autor{Engel} lineal en $\ln\overline{W}$\\
QUAIDS & $E$ & sí & sí & no & más parámetros\\\bottomrule
\end{tabular}\par}
\begin{itemize}
\item \concepto{Sistema completo}: $\sum p_ix_i=\overline{W}$ ($\sum s_i=1$). Agrupar ($\sim n^2$ parámetros): \concepto{separabilidad débil} y presupuestación en dos etapas (\autor{Strotz}, \autor{Gorman}).
\item \concepto{Datos}: microdatos de la EPF (INE); precios: años y regiones, o valores unitarios (\autor{Deaton}).
\item \concepto{Econometría}: covarianzas singulares $\Rightarrow$ se quita una ecuación (MV: invariante, \autor{Barten}); gasto endógeno (instrumento: renta); ceros (censura); demografía (\autor{Pollak}-\autor{Wales}).
\item \concepto{Resultados}: con datos agregados se rechaza sobre todo la \textbf{homogeneidad} (\autor{Barten}, \autor{Deaton} 1974, AIDS 1980); \autor{Deaton}: rigidez, dinámica y \emph{agregación} $\Rightarrow$ microdatos; con microdatos, mejor panorama; contrastes no paramétricos de preferencia revelada (\autor{Blundell} et al., 2003).
\item \concepto{Usos}: IVA (elasticidades compensadas cruzadas), incidencia distributiva, índices por hogar, pobreza. España: AIDS con EPF para reformas del IVA (\autor{Sanz}-\autor{Romero}, 2012).
\end{itemize}

\sub{II.3. Valoración}
\begin{itemize}
\item \textbf{Aportes}: la teoría pasa a ser instrumento de medida; restricciones contrastables; AIDS y QUAIDS concilian flexibilidad, agregación y sencillez.
\item \textbf{Límites}: flexibilidad vs.\ parsimonia ($n(n+1)/2$ términos cruzados; flexibilidad solo local); agregación (forma polar de \autor{Gorman}, PIGLOG; SMD); marco estático (hábitos, duraderos, liquidez).
\end{itemize}

% ==========================================================================
\bloque[12]{III. Aplicaciones de la dualidad}
\sub{III.1. Integrabilidad}
\begin{itemize}
\item ¿Cuándo procede $\vect{x}(\vect{p},\overline{W})$ de maximizar una $U$? \autor{Antonelli} (1886); \autor{Hurwicz}-\autor{Uzawa} (1971).
\item \concepto{Teorema}: $\vect{x}$ $C^1$ es demanda de una $U$ creciente y cuasicóncava $\Leftrightarrow$ (1) ley de \autor{Walras}, (2) $S$ \textbf{simétrica}, (3) $S$ \textbf{semidef.\ negativa} (+ homogeneidad, necesaria). Necesarias y suficientes.
\end{itemize}
\begin{formula}
\[S_{ij}=\frac{\partial x_i}{\partial p_j}+x_j\frac{\partial x_i}{\partial\overline{W}}\quad\text{(observable)}\]
\[\text{EDP: }\nabla_{\vect{p}}E(\vect{p},\overline{U})=\vect{x}\big(\vect{p},E(\vect{p},\overline{U})\big),\quad E(\vect{p}^0,\overline{U})=\overline{W}^0\]
\[U(\vect{x}^0)=\min_{(\vect{p},\overline{W})}V(\vect{p},\overline{W})\ \text{s.a.}\ \vect{p}\cdot\vect{x}^0\leq\overline{W}\]
\end{formula}
\begin{itemize}
\item Demostración (por dualidad): (1) EDP con solución $\Leftrightarrow$ simetría (\concepto{teorema de \autor{Frobenius}}); (2) que sea función de gasto: falta la concavidad, y su hessiana es $S$ $\Rightarrow$ negatividad; (3) invertir $E$ $\to$ $V$ $\to$ $U$.
\end{itemize}
\graficorepaso[0.5\linewidth]{recuperar-utilidad}{Recuperar $U$ desde $V$: el mínimo de $V$ entre los presupuestos con los que se compra $\vect{x}^0$}
\begin{itemize}
\item Con 2 bienes, simetría automática (\autor{Walras} + homogeneidad): solo restringe con $\geq3$.
\item Preferencia revelada: ADPR $\Rightarrow$ $S$ semidef.\ neg., \emph{no} simétrica; AFPR (\autor{Houthakker}) sí; \autor{Kihlstrom}-\autor{Mas-Colell}-\autor{Sonnenschein} los unen; \autor{Afriat} (datos finitos).
\item Relevancia: cierra el círculo (teórica); contraste sin hallar $U$ (empírica); bienestar exacto (\autor{Hausman}, 1981) (normativa).
\end{itemize}

\sub{III.2. Ecuación de \autor{Slutsky}}
\begin{itemize}
\item Derivar $h_i(\vect{p},\overline{U})\equiv x_i\big(\vect{p},E(\vect{p},\overline{U})\big)$ respecto a $p_j$; evaluar en $\overline{U}=V$ y usar \autor{Shephard} ($\partial E/\partial p_j=h_j=x_j$):
\end{itemize}
\begin{formula}
\[\underbrace{\frac{\partial x_i}{\partial p_j}}_{\text{ET}}=\underbrace{\frac{\partial h_i}{\partial p_j}}_{\text{ES}}\ \underbrace{-\,x_j\frac{\partial x_i}{\partial\overline{W}}}_{\text{ER}}\]
\[\varepsilon_{x_i,p_j}=\varepsilon^h_{i,p_j}-s_j\,\varepsilon_{x_i,\overline{W}}\]
\end{formula}
\begin{itemize}
\item ES inobservable; ET y ER observables $\Rightarrow$ $S$ calculable. El ER pesa más cuanto mayor es $s_j$. Compensación de \autor{Hicks} (utilidad) $=$ de \autor{Slutsky} (poder adquisitivo) en el margen.
\item \textbf{Propio}: ES $\leq0$; \textcolor{tceeverde}{\textbf{normal}} $\Rightarrow$ \textcolor{tceeazul}{\textbf{ordinario}} (ley de la demanda); \textcolor{tceedorado}{\textbf{inferior}}: ordinario si domina ES, \textcolor{tceerojo}{\textbf{\autor{Giffen}}} si domina ER ($s_i$ alto). \autor{Giffen} $\Rightarrow$ inferior, no al revés. \autor{Veblen} fuera (precio en las preferencias).
\end{itemize}
\par\vspace{2pt}{\centering\renewcommand{\arraystretch}{1.15}\setlength{\tabcolsep}{3pt}%
\begin{tabular}{@{}ccl@{}}
\toprule
\textbf{ES propio} & \textbf{ER propio} & \textbf{ET propio}\\\midrule
$\leq0$ & \textcolor{tceeverde}{$<0$: normal} & \textcolor{tceeazul}{$\leq0$: ordinario}\\
$\leq0$ & $=0$: frontera & \textcolor{tceeazul}{$\leq0$: ordinario}\\
$\leq0$ & \textcolor{tceedorado}{$>0$: inferior} & \textcolor{tceeazul}{ordinario} si |ES|$\,\geq\,$ER\\
 & & \textcolor{tceerojo}{$>0$: \autor{Giffen}} si ER$\,>\,$|ES|\\\bottomrule
\end{tabular}\par}
\begin{itemize}
\item \textbf{Cruzado}: neto (signo de $S_{ij}$) \textbf{simétrico} (\autor{Hicks}-\autor{Allen}, 1934); bruto no (ER cruzados distintos). \textcolor{tceeverde}{Sustitutivos netos} + $i$ inferior $\Rightarrow$ \textcolor{tceeverde}{sustitutivos brutos}; \textcolor{tceedorado}{complementarios netos} + $i$ normal $\Rightarrow$ \textcolor{tceedorado}{complementarios brutos}; sustitutivos netos + $i$ normal, o complementarios netos + $i$ inferior: ambiguo. Puede haber sustitutivos netos y complementarios brutos.
\item $S\vect{p}=\vect{0}\Rightarrow\sum_{j\neq i}p_jS_{ij}=-p_iS_{ii}\geq0$: todo bien tiene un \concepto{sustitutivo neto}; con 2 bienes, sustitutivos netos.
\item Pendientes en $(x_i,p_i)$: bien normal, marshalliana más plana (más elástica) que la hicksiana; inferior, más inclinada; sin ER, coinciden.
\end{itemize}

\sub{III.3. Bienestar}
\begin{itemize}
\item Un consumidor y cambios de precios (caso particular del equilibrio parcial). Utilidad ordinal: $\overline{U}^1-\overline{U}^0$ no significa nada $\Rightarrow$ medida en euros.
\item Notación: $\vect{p}^0\to\vect{p}^1$, $\overline{W}$ fija; $\overline{U}^k=V(\vect{p}^k,\overline{W})$; cestas $\vect{x}^0$, $\vect{x}^1$.
\end{itemize}
\pareja{bienestar-problema}{Baja $p_1$: de $\overline{U}^0$ a $\overline{U}^1$, pero ¿cuánto?}{bienestar-delta-ec}{$\Delta$EC: área a la izquierda de la marshalliana}
\begin{itemize}
\item \concepto{$\Delta$EC} (\autor{Dupuit}, \autor{Marshall}): $\displaystyle\int_{p_1^1}^{p_1^0}x_1(p_1,\overline{\vect{p}}_{-1},\overline{W})\,dp_1$. Observable, pero mezcla ES y ER; exacta si el ER es nulo (cuasilineal); con varios precios, depende del camino (salvo homotéticas o ER nulos).
\item \concepto{Métrica monetaria} (\autor{McKenzie}, \autor{Samuelson}): $m(\bar{\vect{p}};\vect{p},\overline{W})\equiv E\big(\bar{\vect{p}},V(\vect{p},\overline{W})\big)$: transformación creciente de $V$ en euros. VC: referencia $\vect{p}^1$; VE: referencia $\vect{p}^0$.
\end{itemize}
\begin{formula}
\[\mathrm{VC}=\overline{W}-E(\vect{p}^1,\overline{U}^0)=\int_{p_1^1}^{p_1^0}h_1(p_1,\overline{\vect{p}}_{-1},\overline{U}^0)\,dp_1\]
\[\mathrm{VE}=E(\vect{p}^0,\overline{U}^1)-\overline{W}=\int_{p_1^1}^{p_1^0}h_1(p_1,\overline{\vect{p}}_{-1},\overline{U}^1)\,dp_1\]
\end{formula}
\begin{itemize}
\item \concepto{VC}: precios \emph{finales}, utilidad \emph{inicial}; si baja $p$, lo máximo que pagaría; si sube, $-$VC es la compensación mínima. Área bajo la hicksiana \emph{inicial}.
\item \concepto{VE}: precios \emph{iniciales}, utilidad \emph{final}; si baja $p$, lo mínimo que aceptaría por renunciar; si sube, $-$VE es lo máximo que pagaría por evitarla. Área bajo la hicksiana \emph{final}.
\end{itemize}
\graficorepaso[0.7\linewidth]{bienestar-vc-ve}{VC, $\Delta$EC y VE (y $L$, $P$) ante una bajada de $p_1$ de un bien normal}
\begin{formula}
\[L\leq\mathrm{VC}\leq\Delta\mathrm{EC}\leq\mathrm{VE}\leq P\quad\text{(bien normal, baja }p_1\text{)}\]
\end{formula}
\begin{itemize}
\item Inferior: el orden se invierte; si sube el precio, se invierte en valor absoluto; ER nulo: VC $=\Delta$EC $=$ VE.
\item \concepto{\autor{Willig}} (1976): error relativo de usar $\Delta$EC $\approx|\varepsilon_{x_1,\overline{W}}|\,|\Delta\mathrm{EC}|/2\overline{W}$ (5\,\% de la renta y $\varepsilon=1$: 2,5\,\%):
\end{itemize}
\begin{formula}
\[\mathrm{VC}\approx\Delta\mathrm{EC}-\frac{\varepsilon_{x_1,\overline{W}}(\Delta\mathrm{EC})^2}{2\overline{W}},\qquad \mathrm{VE}\approx\Delta\mathrm{EC}+\frac{\varepsilon_{x_1,\overline{W}}(\Delta\mathrm{EC})^2}{2\overline{W}}\]
\end{formula}
\begin{itemize}
\item Justifica el excedente en equilibrio parcial y coste-beneficio. \autor{Hausman} (1981): falla en la pérdida irrecuperable $\Rightarrow$ $E$ exacta resolviendo la EDP.
\item \textbf{¿Cuál?} Ordenar políticas: \textbf{VE} (base fija $\vect{p}^0$; la VC puede ordenar mal). Compensar (\autor{Kaldor}-\autor{Hicks}): \textbf{VC} (precios que regirán). Exceso de gravamen: $\mathrm{EG}=-\mathrm{VE}-T>0$.
\item \concepto{Info.\ incompleta} (solo $\vect{x}^0$, $\vect{x}^1$): $L=(\vect{p}^0-\vect{p}^1)\cdot\vect{x}^0\leq\mathrm{VC}$, $\mathrm{VE}\leq P=(\vect{p}^0-\vect{p}^1)\cdot\vect{x}^1$, para cualquier bien: \emph{sesgo de sustitución}.
\end{itemize}
\begin{formula}
\[I_K(\overline{U})=\frac{E(\vect{p}^1,\overline{U})}{E(\vect{p}^0,\overline{U})}\ \ \text{(\autor{Konüs}, 1924)}\]
\[I_L=\frac{\vect{p}^1\cdot\vect{x}^0}{\vect{p}^0\cdot\vect{x}^0}\geq I_K(\overline{U}^0),\qquad I_P=\frac{\vect{p}^1\cdot\vect{x}^1}{\vect{p}^0\cdot\vect{x}^1}\leq I_K(\overline{U}^1)\]
\end{formula}
\graficorepaso[0.72\linewidth]{indice-konus-laspeyres}{\autor{Laspeyres} sobrestima: el sesgo de sustitución}
\begin{itemize}
\item $I_K(\overline{U}^0)$: versión de la VC; $I_K(\overline{U}^1)$: de la VE. \autor{Laspeyres}: aprox.\ de 1.er orden de $E(\vect{p}^1,\overline{U}^0)$ $\Rightarrow$ cota superior (concavidad). Solo con homotéticas $I_P\leq I_K\leq I_L$; igualdades solo con \autor{Leontief}.
\item \concepto{Superlativos} (\autor{Diewert}, 1976): \autor{Fisher} $\sqrt{I_LI_P}$ (exacto para $E$ cuadrática homogénea); \autor{Törnqvist} $\ln I_T=\sum\tfrac12(s_i^0+s_i^1)\ln(p_i^1/p_i^0)$ (translog).
\item IPC del INE: \autor{Laspeyres} encadenado (base 2025; ponderaciones anuales: reducen el sesgo, no lo eliminan) $\Rightarrow$ sobrevalora. Deflactor del PIB: \autor{Paasche} encadenado $\Rightarrow$ infravalora.
\end{itemize}
\begin{ideaclave}
$E$ traduce la utilidad a euros: VC y VE exactas; $\Delta$EC las aproxima tanto mejor cuanto menor es el ER (\autor{Willig}); \autor{Laspeyres} y \autor{Paasche} las acotan con un sesgo de sustitución que sigue en los índices oficiales.
\end{ideaclave}

% ==========================================================================
\bloque[1]{Conclusión}
\begin{itemize}
\item \textbf{Recapitulación}: hicksiana, $E$ y relaciones de dualidad (I); sistemas de demanda, AIDS y microdatos (II); integrabilidad, \autor{Slutsky} y bienestar (III).
\item \textbf{Relevancia}: de lo inobservable (utilidad) a restricciones sobre lo observable y bienestar en euros: elasticidades fiscales, exceso de gravamen, coste-beneficio, sesgo del IPC.
\item \textbf{Extensiones}: empresa \vertema{3.A.12}; equilibrio parcial \vertema{3.A.16}; bienestar \vertema{3.A.22}\vertema{3.A.24}; coste-beneficio \vertema{4.B.6}; exceso de gravamen \vertema{4.B.9}; IPC y deflactor \vertema{3.A.30}\vertema{4.A.1}; riesgo \vertema{3.A.10}; intertemporal \vertema{3.A.29}.
\item \textbf{Opinión}: teoría abstracta muy útil; límite: consumidor estático y agregación (sumar exige juicios de valor); herencia: \autor{Deaton} (microdatos, heterogeneidad).
\item \textbf{Idea final}: la pensión con IPC sobrecompensa algo si la cesta se parece a la del IPC (\autor{Laspeyres}); si es distinta, puede quedarse corta. La compensación exacta es la VC.
\end{itemize}

% ==========================================================================
\bloque{Autores y fechas}
\hito{1844}{\autor{Dupuit}: excedente del consumidor}
\hito{1915}{\autor{Slutsky}: ecuación ES + ER}
\hito{1924}{\autor{Konüs}: índice verdadero de coste de la vida}
\hito{1939}{\autor{Hicks}: demanda compensada; VC y VE (1941-43)}
\hito{1947}{\autor{Roy}: identidad de Roy}
\hito{1953}{\autor{Shephard}: lema; 1954, \autor{Stone}: LES}
\hito{1957}{\autor{McKenzie}: función de gasto}
\hito{1971}{\autor{Hurwicz}-\autor{Uzawa}: integrabilidad}
\hito{1976}{\autor{Willig}: aproximación; \autor{Diewert}: superlativos}
\hito{1980}{\autor{Deaton}-\autor{Muellbauer}: AIDS; 1981, \autor{Hausman}: bienestar exacto}
\hito{1996}{Comisión \autor{Boskin}: sesgo del IPC}
\hito{2015}{Nobel a \autor{Deaton}}

\bloque{Datos clave}
\begin{itemize}
\item Revalorización de las pensiones: \dato{2,7\,\% en 2026}{media del IPC, LGSS}.
\item EPF: \dato{unos 24.000 hogares}{INE}, anual desde 2006; gasto medio \dato{35.101 €/hogar}{INE, EPF 2025}.
\item IPC base 2025: \dato{487 artículos, 196 subclases, 13 grupos}{INE, 2026}.
\item Deflactor del PIB: \dato{+3,0\,\%}{INE, 2025}.
\item Sesgo del IPC de EE.\,UU.: \dato{1,1 p.p./año, 0,4 de sustitución}{\autor{Boskin}, 1996}; \dato{0,18 p.p./año}{\autor{Manser}-\autor{McDonald}, 1959-85}; IPC encadenado \autor{Törnqvist} \dato{0,25 p.p. menos/año}{BLS, 2005-24}.
\end{itemize}

\begin{trampa}
\begin{itemize}
\item $E$: homog.\ de grado \emph{1} en $\vect{p}$; $\vect{h}$: grado \emph{0} en $\vect{p}$; $V$ y $\vect{x}$: grado 0 en $(\vect{p},\overline{W})$.
\item \autor{Shephard} sin signo menos; \autor{Roy}, con signo menos.
\item $S$: \emph{primeras} derivadas de $\vect{h}$, \emph{segundas} de $E$.
\item $V$ es \emph{cuasiconvexa} (no cuasicóncava) y no depende de la convexidad de las preferencias; $E$ es \emph{cóncava} en $\vect{p}$.
\item ER nulo: VC $=\Delta$EC $=$ VE, pero $L\neq P$ si hay sustitución.
\item \autor{Laspeyres} y \autor{Paasche} acotan índices verdaderos \emph{distintos}; $I_P\leq I_K\leq I_L$ solo con homotéticas.
\item Con datos agregados se rechaza sobre todo la \emph{homogeneidad}, no la simetría.
\item $\lambda^*=\partial V/\partial\overline{W}$; $\mu^*=\partial E/\partial\overline{U}$; $\mu^*=1/\lambda^*$.
\item ADPR $\Rightarrow$ $S$ semidef.\ neg., \emph{no} simétrica.
\item Teorema de \autor{Frobenius}, no de \autor{Rouché}-\autor{Frobenius}.
\item El LES no admite bienes inferiores ($\alpha_i>0$) ni complementarios netos.
\item \autor{Giffen} $\Rightarrow$ inferior; inferior $\not\Rightarrow$ \autor{Giffen}.
\end{itemize}
\end{trampa}

% ==========================================================================
\bloque{Preguntas del tribunal}
\begin{enumerate}
\item \textbf{¿Por qué $\mu^*=1/\lambda^*$?} Derivando $E(\vect{p},V(\vect{p},\overline{W}))=\overline{W}$ respecto a $\overline{W}$: $\frac{\partial E}{\partial\overline{U}}\frac{\partial V}{\partial\overline{W}}=1$.
\item \textbf{¿Por qué la hicksiana no admite \autor{Giffen}?} No tiene efecto renta y $E$ es cóncava: $\partial h_i/\partial p_i\leq0$.
\item \textbf{¿Por qué $S$ es simétrica?} Son segundas derivadas de $E$ (\autor{Schwarz}); las brutas no, porque el ER cruzado difiere.
\item \textbf{¿Qué condiciones da la integrabilidad?} Ley de \autor{Walras}, simetría y negatividad de $S$; con 2 bienes la simetría es automática.
\item \textbf{¿Por qué es «casi ideal» el AIDS?} Flexible, coherente, agrega exactamente y es lineal (LA-AIDS), pero la negatividad no está garantizada y las curvas de \autor{Engel} son restrictivas.
\item \textbf{¿Cómo se demuestra la ecuación de \autor{Slutsky} por dualidad?} Derivando $h_i(\vect{p},\overline{U})\equiv x_i(\vect{p},E(\vect{p},\overline{U}))$ respecto a $p_j$ y usando \autor{Shephard}.
\item \textbf{¿VC o VE?} VE para ordenar políticas (base fija); VC para compensar a los perjudicados.
\item \textbf{¿Cómo se mide el exceso de gravamen?} Con la VE: $\mathrm{EG}=-\mathrm{VE}-T>0$; el consumidor pagaría más de $T$ por un impuesto de suma fija equivalente.
\item \textbf{¿Cuándo vale el excedente marshalliano?} Si el ER es pequeño (\autor{Willig}): bien que pesa poco en la renta.
\item \textbf{¿Por qué se rechazaba la homogeneidad con datos agregados?} Según \autor{Deaton}, por especificaciones rígidas, la dinámica y sobre todo la agregación: de ahí los microdatos.
\item \textbf{¿Por qué el IPC sobrestima el coste de la vida?} \autor{Laspeyres} fija la cesta e ignora la sustitución; la concavidad de $E$ lo hace cota superior.
\end{enumerate}


\end{columnas}
\end{document}
